\documentclass[10pt,a4paper]{amsart}
\usepackage[margin=19mm]{geometry}
\usepackage{amsmath,amssymb,mathtools}
\usepackage[scr=rsfs,cal=euler]{mathalfa}
\usepackage{xfrac}
\usepackage[unicode,hidelinks,bookmarks=false]{hyperref}
\newcommand{\ie}{i.e.\ }
\newcommand{\cf}{cf.\ }
\newcommand{\resp}{respectively\ }
\newcommand{\Subsection}{\subsection}
\providecommand{\nobreakdash}{\nobreak\hbox{-}}
\setlength{\emergencystretch}{2em}
\multlinegap0pt
\usepackage{bm}
\usepackage{bbm}
\usepackage{dsfont}
\usepackage{xspace}
\usepackage{cancel}
\usepackage{mathtools}
\usepackage{amsthm}
\makeatletter
\@ifundefined{theorem}{\newtheorem{theorem}{Theorem}}{}
\@ifundefined{lemma}{\newtheorem{lemma}[theorem]{Lemma}}{}
\makeatother
\newcommand{\E}{\mathbb{E}}
\newcommand{\PP}{\mathbb{P}}
\newcommand{\V}{\mathbb{V}}
\newcommand{\prob}[1]{\mathbb{P}\left( #1 \right)}
\newcommand{\procond}[2]{\mathbb{P}\mathopen{}\left[#1\middle|#2\right]}
\title[Subtractive random forests with two choices]{Subtractive random forests with two choices}
\author{Francisco Calvillo, Luc Devroye, G\'abor Lugosi}
\date{Source: arXiv:2405.10455v2 (CC BY 4.0)}
\begin{document}
\maketitle

\noindent\emph{Source and licence.} Subtractive random forests with two choices. Version arXiv:2405.10455v2 (CC BY 4.0), distributed under the Creative Commons Attribution 4.0 International licence, as recorded in the arXiv licence metadata for this exact version and shown on the abstract page https://arxiv.org/abs/2405.10455v2. Full rights notice: licenses/arxiv-2405.10455-CC-BY-4.0.txt.\par\smallskip

\noindent\textit{Editorial scope.} Counted unit: the Lemma whose source label is \texttt{Lemma second moment} in the article (source0.tex lines 781--789, source label \texttt{Lemma second moment}), delivered together with the article's own complete original proof of it (source0.tex lines 791--873, one \texttt{proof} environment of 83 lines). No mathematics is written, repaired or re-derived: statement and proof are the article's own text line for line. Required context retained uncounted: convergence vn (lines 767--769); lowerbound (lines 771--773). Every definition, display and label of those lines is retained unchanged. Omitted from this package and not redistributed: 1--766, 770--770, 774--780, 790--790, 874--1279. Nothing inside a retained excerpt is omitted or altered: every retained body is delivered exactly as the article's lines are written, so no omission receipt of any kind accompanies this package. Citations: the retained excerpts carry no citation command at all, so no bibliography entry of the article is transcribed and no reference is redistributed. Reference adaptation: none was needed and none was made; every reference inside the delivered unit resolves inside it, so no unresolved reference or citation is produced. Printed slips kept literal: no printed word, symbol, hypothesis or number of the counted statement or of its proof was altered. Layout and class: the article's own class is \texttt{article} with packages including kpfonts, titlesec, dsfont, comment, textcomp, amssymb, bbm, fancyhdr, amsmath, amsbsy, amsthm, amscd, latexsym, graphicx; the wrapper is the lane's pinned \texttt{amsart} editorial wrapper at 10pt on A4, which restates the theorem-like environments the retained text needs and adds the article's own macro definitions that the retained text needs (\texttt{\textbackslash E}, \texttt{\textbackslash PP}, \texttt{\textbackslash V}, \texttt{\textbackslash prob}, \texttt{\textbackslash procond}), with extra wrapper packages (bm, bbm, dsfont, xspace, cancel, mathtools). The delivered numbering is the unit's own. Assets: no retained span contains a figure environment, a graphics-inclusion command, a PDF-page inclusion, a table or a TikZ picture; no image file is copied into this package, the package declares no asset dependency and no image file is redistributed with it. Licence: version arXiv:2405.10455v2 of this article is distributed under the Creative Commons Attribution 4.0 International licence, as recorded in the arXiv licence metadata for this exact version and shown on the abstract page https://arxiv.org/abs/2405.10455v2. A full rights notice is delivered at licenses/arxiv-2405.10455-CC-BY-4.0.txt. Characters: every retained excerpt is pure ASCII, so the delivered engine, XeLaTeX with its default Latin Modern text font, reports no missing character.
\begin{equation}
\lim_{n \to \infty} v_n =   \frac{1}{\E \min (Z,W)} >0~. \label{convergence vn}
\end{equation}

\begin{equation}
\E J_n \ge c \sum_{m=1}^n p_m \to \infty \quad \text{as } n\to \infty~. \label{lowerbound}
\end{equation}

\begin{lemma}\label{Lemma second moment}
Assume that $\E Z = \infty$ and $\E \min (Z,W) < \infty$, and let $\V J_n$ denote the variance of the random variable $J_n$. Then
\begin{equation}
\frac{\V J_n}{(\E J_n)^2}\to 0 \quad \text{as } n \to \infty~. \label{scnd moment method}
\end{equation}
In particular, we have:
$$\PP\left(J_n\le \frac{\E J_n}{2}\right)=\PP\left(J_n - \E J_n \le - \frac{\E J_n}{2}\right)\le 4\frac{\V J_n}{(\E J_n)^2}\to 0~,$$
and therefore $J_n \to \infty$ in probability as $n \to \infty$.
\end{lemma}

\begin{proof}[Proof of Lemma \ref{Lemma second moment}]
To show \eqref{scnd moment method}, we prove that $\E[J_n^2]/(\E J_n)^2\to 1$ as $n \to \infty$. We first define
$$
p_n^*=\PP(\max(Z,W)\ge n) = 2p_n - p_n^2~.
$$
Recall that $v_n = \PP(0 \in S_n)$, as introduced earlier. Then, for $n\ge 1$,
\begin{align}
    (\E J_n)^2 &= 2\sum_{1\le k<m\le n} p_m^* p_k^* v_{n-m}v_{n-k} + \sum_{1\le m \le n} (p_m^*)^2 v_{n-m}^2 \nonumber \\
    &= 2\sum_{1\le k<m\le n} p_m^* p_k^* v_{n-m}v_{n-k} + O(1)
\end{align}
since $\sum_{1\le m \le n} (p_m^*)^2 v_{n-m}^2$ is bounded by $4\E
\min (Z,W) <\infty$.

Now, define for any $m\ge 1$ the random variable $X_m=\max(Z_m, W_m)$ and observe that for any $1\le k < m \le n$, the events $\left\{ X_k\ge k\right\}$ and $\left\{ X_m\ge m, m\in S_n, k \in S_n\right\}$ are independent.
Thus, for any fixed $n\ge 1$, we may write
\begin{align}
    \E\left[ J_n^2\right] &= 2\sum_{1\le k < m \le n} \prob{m\in S_n, k\in S_n, X_m\ge m, X_k\ge k} + \E J_n \nonumber \\
    &= 2\sum_{1\le k < m \le n} p_k^*\prob{m\in S_n, k\in S_n, X_m\ge m} + \E J_n \nonumber \\
    &= 2\sum_{1\le k < m \le n} p_k^* \procond{k\in S_n, X_m \ge m}{m\in S_n}\prob{m\in S_n}+ \E J_n \nonumber 
\end{align}
If $m\in S_n$, the event $\{k\in S_n, X_m\ge m\}$ can only happen if one of the two random variables $Z_m$ and $W_m$ is greater or equal to $m$ and the other one is smaller than $m-k$, for otherwise $k$ could not belong to $S_n$ when $1\le k<m \le n$. Thus,
using the independence of $Z_m$ with respect to $W_m$, $\{k\in S_n\}$ and $\{m\in S_n\}$, and using the fact that 
$p_m^* + p_m^2 = 2p_m$,
\begin{align}
     \E&\left[ J_n^2\right] \\
     &= 2\sum_{1\le k < m \le n} p_k^* \times 2 \procond{Z_m\ge m, W_m\le m-k , k\in S_n}{m\in S_n}\prob{m\in S_n}+ \E J_n \nonumber \\
     &= 2\sum_{1\le k < m \le n} p_k^* (p_m^*+p_m^2)v_{n-m}\procond{W_m\le m-k , k\in S_n}{m\in S_n}+ \E J_n \nonumber \\
    &= 4\sum_{1\le k < m \le n} p_k^* (p_m^*+p_m^2)v_{n-m}\sum_{i=1}^{k-m}\prob{k\in S_{m-i}}\prob{W_m=i}+ \E J_n \nonumber \\
    &= 2\sum_{1\le k < m \le n} p_k^* (p_m^*+p_m^2)v_{n-m}\sum_{i=1}^{m-k}v_{m-k - i}q_i+ \E J_n ~.
\end{align}
Observe that
\begin{align}
    \sum_{1 \le k < m \le n}p_k^*p_m^2v_{n-m}\sum_{i=1}^{m-k}q_iv_{m-k-i} &\le \sum_{1 \le k < m \le n}p_k^*p_m^2v_{n-m} \nonumber \\
    &\le  \sum_{1 \le k < m \le n}p_k^*p_m^2 \nonumber \\
    &\le \left( \sum_{m=1}^\infty p_m ^2\right) \left(\sum_{k=1}^n p_k^*\right) \nonumber \\
    &=  O\left( \E J_n\right) \nonumber
\end{align}     
by using the inequality in \eqref{lowerbound}, which leads to 
\begin{equation}
    \E[J_n^2] = 2\sum_{1\le k < m \le n} p_k^*p_m^*v_{n-m}\sum_{i=1}^{m-k}v_{m-k-i}q_i+ O\left( \E J_n \right)~.
\end{equation}
Fix some $\epsilon >0$. By \eqref{convergence vn}, we know that there exists some constant $x>0$ such that for all $m \ge x$, we have $|v_m-\lambda|<\epsilon$,  where $\lambda = 1/(\E\min(Z,W))$.

\noindent If $m\ge x$, see that
\begin{align*}
    \sum_{i=1}^m q_i v_{m-i} &\le \sum_{i=1}^{m-x} q_i (\lambda + \epsilon) + \sum_{m-x<i\le m} q_i \\
    &\le (\lambda + \epsilon) + p_{m-x}~.
\end{align*} 
Thus, there exists some constant $y>x$ such that for all $m \ge y$,
$$ \sum_{i=1}^m q_i v_{m-i}\le \lambda + 2 \epsilon~.$$
Moreover,
\begin{align*}
    \sum_{\substack{1\le k < m \le n \\ \textit{s.t. } m-k < y} }p_m^*v_{n-m}\underbrace{p_k^*\sum_{i=1}^{m-k}v_{m-k-i}q_i}_{\le 1}  
    &\le \sum_{1\le m \le n}p_m^* v_{n-m} \sum_{k=m-y}^m 1 \\
    &= (y+1) \E J_n = O\left( \E J_n\right)~,
\end{align*}
and 
\begin{align*}
    \sum_{\substack{1\le k < m \le n \\ \textit{s.t. } n-m < y} }p_k^*\underbrace{p_m^*v_{n-m}\sum_{i=1}^{m-k}v_{m-k-i}q_i}_{\le 1} &\le \sum_{n-y\le m \le n}1 \sum_{k=1}^n p_k^* \\
    &\le (y+1) \sum_{k=1}^n p_k^*
    = O\left( \E J_n\right)~.
\end{align*}
Putting everything together, we have that 
\begin{align}
    \E [J_n^2] = \sum_{m=1}^{n-y}\sum_{k=1}^{m-y} p_m^*p_k^*v_{n-k}\sum_{i=1}^{m-k} q_i v_{m-k-i} + O\left( \E J_n\right)~. \label{equivalent 1}
\end{align}
Similarly, we have
\begin{align}
    (\E J_n)^2 = \sum_{s=1}^{n-y}\sum_{k=1}^{m-y} p_m^*p_k^*v_{n-m}v_{n-k} + O\left( \E J_n\right)~. \label{equivalent 2}
\end{align}

\noindent Finally, observe that for any $n> y$ we have:
\begin{align}
\sum_{m=1}^{n-y}\sum_{k=1}^{m-y} p_m^*p_k^*\underbrace{v_{n-m}\sum_{i=1}^{m-k} q_i v_{m-k-i}}_{\le (\lambda + 2\epsilon)^2} \le (\lambda +2\epsilon)^2\sum_{m=1}^{n-y}\sum_{k=1}^{m-y} p_m^*p_k^* \label{upper J squared}
\end{align}
and
\begin{align}
    \sum_{m=1}^{n-y}\sum_{k=1}^{m-y} p_m^*p_k^*\underbrace{v_{n-m}v_{n-k}}_{\ge (\lambda - \epsilon)^2} \ge (\lambda - \epsilon)^2\sum_{m=1}^{n-y}\sum_{k=1}^{m-y} p_m^*p_k^*~. \label{lower J squared}
\end{align}
From \eqref{equivalent 1}, \eqref{equivalent 2}, \eqref{upper J squared} and \eqref{lower J squared} it follows that
$$\limsup_{n\to \infty} \frac{\E[J_n^2]}{(\E J_n)^2}\le \left( \frac{\lambda + 2\epsilon}{\lambda - \epsilon}\right)^2 ~.$$
Since this is true for any $\epsilon>0$ small enough, and $\E[J_n^2]\ge (\E J_n)^2$, we conclude that $$\frac{\E [J_n^2]}{(\E J_n)^2}\to 1$$
as $n \to \infty$.\end{proof}

\end{document}
